\chapter{总结与展望}\label{chap:conclusion}

\section{总结}

语音到语音生成旨在实现从源语音到目标语音的直接转换，涵盖语音到语音翻译和语音交互等典型任务，是构建自然人机交互系统的核心技术。随着模型结构从编码器-解码器架构演进到以大语言模型为核心的仅解码器架构，语音到语音生成在翻译质量、响应延迟、语音自然度和交互灵活性等方面不断取得突破，但仍面临多个尚未解决的关键问题。

本文围绕语音到语音生成的模型结构演进，从编码器-解码器架构下的语音到语音翻译、仅解码器架构下的实时语音交互、以及面向全双工交互的语音大模型三个方向展开系统性研究，共完成了五项具体工作。

\textbf{1. 基于词表转换的可组合语音到语音翻译框架}

端到端语音到语音翻译模型需要稀缺的平行语音翻译数据从头训练，无法复用语音到文本翻译和语音合成领域已有的预训练模型，根本原因在于不同模型使用不同的词表和表示空间。针对这一问题，本文第\ref{chap:comspeech}章提出了ComSpeech，通过基于CTC的词表适配器实现不同词表之间的表示转换，使任意的语音到文本翻译模型和语音合成模型可以自由组合。在零样本场景下，ComSpeech-ZS仅使用各自领域的数据即可构建完整的语音到语音翻译系统。在CVSS数据集上，ComSpeech在有监督学习场景下超越了所有基线模型，ComSpeech-ZS在零样本场景下仍超越了所有有监督基线。

\textbf{2. 基于有向无环图的语音到语音翻译模型}

在编码器-解码器架构下，非自回归语音到语音翻译方法虽然解码速度快，但由于目标序列具有语言和声学两个层面的多峰分布特性，非自回归模型难以有效建模这种多样性，导致翻译质量与自回归模型存在较大差距。针对这一问题，本文第\ref{chap:daspeech}章提出了DASpeech，采用有向无环图Transformer建模语言层面的多峰分布，通过FastSpeech 2建模声学层面的多峰分布，并提出期望路径训练算法实现端到端训练。在CVSS法语到英语数据集上，DASpeech在翻译质量上达到与自回归模型Translatotron 2相当甚至更优的水平，同时解码速度最高提升18.53倍。

\textbf{3. 基于非自回归流式解码的实时语音交互大模型}

随着大语言模型的兴起，语音交互成为用户与大语言模型之间最自然的交互方式。然而，实现低延迟的实时语音交互面临一个核心难题：大语言模型以自回归方式逐个生成文本词元，若等待文本生成完毕再合成语音，将导致不可接受的响应延迟。针对这一问题，本文第\ref{chap:llama-omni}章提出了LLaMA-Omni，设计了基于CTC的非自回归流式语音解码器，在大语言模型生成每个文本词元时同步预测对应的离散语音单元，实现了文本响应与语音响应的同步流式生成。在InstructS2S-Eval上，LLaMA-Omni在离线场景下的ChatGPT评分达到3.99，显著优于SpeechGPT、SALMONN和Qwen2-Audio等基线系统，流式场景下响应延迟低至236毫秒。

\textbf{4. 基于自回归流式解码的高质量语音交互大模型}

LLaMA-Omni的非自回归语音解码器虽然延迟极低，但由于各位置的预测相互独立，缺乏对语音词元序列依赖关系的建模，导致生成语音的准确性和自然度不足。针对这一问题，本文第\ref{chap:llama-omni2}章提出了LLaMA-Omni 2，设计了基于读写交替策略的自回归流式语音解码器，通过门控融合机制自适应地利用大语言模型的隐状态和文本词元嵌入，在大语言模型生成文本的过程中交替生成对应的语音词元。LLaMA-Omni 2在语音问答和语音指令遵循任务上均超越了LLaMA-Omni和GLM-4-Voice，将语音到文本与语音到语音之间的性能差距大幅缩小，响应延迟低于600毫秒，且模型性能随大语言模型规模的增大而持续提升。

\textbf{5. 基于多信道交织的全双工语音大模型}

上述语音交互模型均面向基于轮次的对话场景，依赖外部的语音活动检测模块判断用户发言边界，无法支持打断、停顿等灵活交互模式。针对这一问题，本文第\ref{chap:full-duplex}章提出了LLaMA-Omni-Duplex，一种基于多信道交织的全双工语音交互模型，将用户语音、模型回复文本和模型回复语音三个信道按固定比例组织成交织序列，统一由大语言模型进行自回归建模。通过文本信道中的特殊标记将全双工场景下的时机决策转化为下一词元预测问题，无需引入额外的分类模块。基于GLM-4-Voice进行轻量的全双工微调，并结合直接偏好优化进一步强化时机决策能力，LLaMA-Omni-Duplex的轮替成功率达到92.0\%，打断成功率达到100.0\%，回复质量显著超越Moshi基线。


\section{展望}

尽管本文在语音到语音生成领域取得了上述进展，但仍存在多个值得深入探索的方向。

\textbf{1. 语音生成的表现力与可控性}

本文提出的语音交互模型生成的语音在音色和风格上较为单一，缺乏对副语言信息的细粒度控制。在实际应用中，用户期望模型能够根据对话场景自适应地调整语气、语速和情感表达，而非始终以固定的风格进行回复。未来的研究可以探索如何在语音交互模型中引入情感识别和风格控制机制，使模型能够感知用户的情绪状态并生成匹配的语音回复，提升交互体验的自然度和亲和力。

\textbf{2. 全双工交互的场景泛化}

本文的全双工模型主要在轮替和打断两类场景上进行了验证，但真实对话中还存在反馈应答、思考停顿、重叠发言等多种复杂现象。虽然本文提出的多信道交织建模方式在理论上兼容这些场景，但由于缺乏相应的训练数据，模型在这些场景上的表现尚待验证。未来的研究可以构建覆盖更多全双工场景的训练数据，并探索如何在有限的全双工数据条件下提升模型的场景泛化能力。

\textbf{3. 多模态信息的融合与理解}

本文的工作聚焦于语音和文本两种模态的交互，但实际应用场景中往往需要同时处理视觉等更多模态的信息。例如，智能助手在语音交互的同时可能需要理解用户分享的图片或视频内容，并据此生成语音回复。如何将视觉等模态高效地融入现有的语音交互框架，在保持低延迟的同时实现跨模态的理解和生成，是构建通用多模态交互系统的重要方向。

\textbf{4. 模型的轻量化与部署效率}

本文提出的模型在参数量和计算开销上仍较大，难以直接部署在移动设备或边缘计算平台上。虽然LLaMA-Omni 2通过多规模模型的设计展示了小参数量模型的可行性，但在设备端实现实时语音交互仍需在模型压缩、推理加速等方面进行更多探索。量化、剪枝、知识蒸馏等模型压缩技术以及针对语音交互场景的推理优化策略，都是值得深入研究的方向。
